% Part: first-order-logic
% Chapter: natural-deduction
% Section: provability-propositional

% verification of properties of provability needed for maximally
% consistent sets in the completeness chapter.

\documentclass[../../../include/open-logic-section]{subfiles}

\begin{document}

\iftag{FOL}
      {\olfileid[hi]{fol}{ntd}{ppr}}
      {\olfileid[hi]{pl}{ntd}{ppr}}

\olsection{\usetoken{S}{derivability} और प्रतिज्ञप्ति-संयोजक}

\begin{explain}
  हम स्थापित करते हैं कि स्वाभाविक निगमन का !!{derivability} संबंध~$\Proves$
  इतना प्रबल है कि प्रतिज्ञप्ति-संयोजकों से जुड़े कुछ आधारभूत तथ्य सिद्ध कर
  सके, जैसे $!A \land !B
  \Proves !A$ और $!A, !A \lif !B \Proves !B$ (पूर्वपक्ष-स्थापन)। पूर्णता
  प्रमेय के प्रमाण के लिए ये तथ्य आवश्यक हैं।
\end{explain}

\begin{prop}\ollabel{prop:provability-land}
  \begin{enumerate}
  \item \ollabel{prop:provability-land-left} दोनों $!A \land !B \Proves
    !A$ और $!A \land !B \Proves !B$ सत्य हैं।
  \item \ollabel{prop:provability-land-right} $!A, !B \Proves !A \land !B$.
  \end{enumerate}
\end{prop}

\begin{proof}
  \begin{enumerate}
  \item दोनों को !!{derive} कर सकते हैं:
    \begin{prooftree}
      \AxiomC{$!A \land !B$}
      \RightLabel{\Elim{\land}}
      \UnaryInfC{$!A$}
      \DisplayProof\qquad\bottomAlignProof
      \AxiomC{$!A \land !B$}
      \RightLabel{\Elim{\land}}
      \UnaryInfC{$!B$}
    \end{prooftree}
  \item यह !!{derive} कर सकते हैं:
    \begin{prooftree}
      \AxiomC{$!A$}
      \AxiomC{$!B$}
      \RightLabel{\Intro{\land}}
      \BinaryInfC{$!A \land !B$}
    \end{prooftree}
  \end{enumerate}
\end{proof}


\begin{prop}\ollabel{prop:provability-lor}
  \begin{enumerate}
  \item $!A \lor !B, \lnot !A, \lnot !B$ असंगत है।
  \item दोनों $!A \Proves !A \lor !B$ और $!B \Proves !A \lor !B$ सत्य हैं।
  \end{enumerate}
\end{prop}

\begin{proof}
  \begin{enumerate}
  \item निम्नलिखित !!{derivation} पर विचार करें:
    \begin{prooftree}
      \AxiomC{$!A \lor !B$}
      \AxiomC{$\lnot !A$}
      \AxiomC{$\Discharge{!A}{1}$}
      \RightLabel{\Elim{\lnot}}
      \BinaryInfC{$\lfalse$}
      \AxiomC{$\lnot !B$}
      \AxiomC{$\Discharge{!B}{1}$}
      \RightLabel{\Elim{\lnot}}
      \BinaryInfC{$\lfalse$}
      \DischargeRule{\Elim{\lor}}{1}
      \TrinaryInfC{$\lfalse$}
    \end{prooftree}
    यह~$\lfalse$ की !!a{derivation} है, !!{undischarged} मान्यताओं
    $!A \lor !B$, $\lnot !A$, और $\lnot !B$ से।
  \item दोनों को !!{derive} कर सकते हैं:
    \begin{prooftree}
      \AxiomC{$!A$}
      \RightLabel{\Intro{\lor}}
      \UnaryInfC{$!A \lor !B$}
      \DisplayProof\qquad\bottomAlignProof
      \AxiomC{$!B$}
      \RightLabel{\Intro{\lor}}
      \UnaryInfC{$!A \lor !B$}
    \end{prooftree}
  \end{enumerate}
\end{proof}

\begin{prop}\ollabel{prop:provability-lif}
  \begin{enumerate}
  \item \ollabel{prop:provability-lif-left}  $!A, !A \lif !B \Proves !B$.
  \item \ollabel{prop:provability-lif-right}
    दोनों $\lnot !A \Proves !A \lif !B$ और $!B \Proves !A \lif !B$ सत्य हैं।
  \end{enumerate}
\end{prop}

\begin{proof}
  \begin{enumerate}
  \item यह !!{derive} कर सकते हैं:
    \begin{prooftree}
      \AxiomC{$!A \lif !B$}
      \AxiomC{$!A$}
      \RightLabel{\Elim{\lif}}
      \BinaryInfC{$!B$}
    \end{prooftree}
    
  \item इसे निम्नलिखित दो !!{derivation}s दिखाती हैं:
    \begin{prooftree}
      \AxiomC{$\lnot !A$}
      \AxiomC{$\Discharge{!A}{1}$}
      \RightLabel{\Elim{\lnot}}
      \BinaryInfC{$\lfalse$}
      \RightLabel{\FalseInt}
      \UnaryInfC{$!B$}
      \DischargeRule{\Intro{\lif}}{1}
      \UnaryInfC{$!A \lif !B$}
      \DisplayProof\qquad\bottomAlignProof
      \AxiomC{$!B$}
      \RightLabel{\Intro{\lif}}
      \UnaryInfC{$!A \lif !B$}
    \end{prooftree}
    ध्यान दें कि $\Intro{\lif}$ मान्यता~$!A$ को !!{discharge} कर सकता है,
    पर ऐसा करना आवश्यक नहीं है।
  \end{enumerate}
\end{proof}

\end{document}
